%%%%%
%
%% Dans ce fichier, les �ventuels accents sont cod�s
%% en ISO-8859-15 (aka latin9)
%% La ligne suivante donne trois exemples 
%% e accent aigu : �, e accent grave : �, e-dans-l'o : �
\section{�chantillonnage (et moyenne) d'une variable \texttt{double} sur une nouvelle base
de temps}

\itemstyle{Op�rateur~:} \operfmt{MoyEch\_tps\_impos\_dbl}
\index{MoyEch\_tps\_impos\_dbl}

%%%%%OBJET%%%%
\itemstyle{Objet~:}
Cette fonction calcule, pour chaque instant $t$ d'une
base de temps "time\_out", 
une variable "X\_out", contenant la moyenne des
valeurs de  "X\_in" entre $t$-�cart\_avant et $t$+�cart\_apr�s.

Aucune supposition n'est faite sur la r�gularit� de
l'�chantillonnage des 2 variables. Cette fonction peut
donc servir, par exemple, � transformer une s�rie temporelle
asynchrone en une s�rie r�guli�re. Il est cependant n�cessaire
que time\_in (le temps de X\_in) soit "� peu pr�s" tri� par ordre
croissant.

%%%%%ARGUMENTS%%%%
\bigskip
%% Dans ce fichier, les �ventuels accents sont cod�s
%% en ISO-8859-15 (aka latin9)
%% La ligne suivante donne trois exemples 
%% e accent aigu : �, e accent grave : �, e-dans-l'o : �
%%%%%ENTREES%%%%
\begin{tabarg}
\multicolumn{3}{l}{\titretab {Entr\'ees}}\\
\hline
$X_{in}$ & Vecteur, \texttt{double} & Variable d'origine\\
$time_{in}$ & Vecteur, \texttt{double} & Base de temps de $X_{in}$ (secondes)\\
$ecart\_avant$ & Coeff., \texttt{double} & �cart maxi avant (secondes)\\
$ecart\_apres$ & Coeff., \texttt{double} & �cart maxi apr�s (secondes)\\
$toler\_mnq$ & Coeff., \texttt{Int\_32} & Seuil de calcul de la moyenne (nb. �chant). A chaque instant de la base de temps d'arriv�e, 
la moyenne est calcul�e si le nombre d'�chantillons dans l'intervalle de temps est strictement sup�rieur � $toler\_mnq$. Sinon, 
le r�sultat est "valeur manquante". Une valeur de 0 provoque le calcul dans tous les cas (m�me s'il n'y a qu'un seul �chantillon).\\
$time_{out}$ & Vecteur, \texttt{double} & Base de temps d'arriv�e (secondes)\\
$size_{out}$ & Coeff., \texttt{Int\_32} & Nombre de valeurs dans la base de temps 
d'arriv�e (dans adc, utiliser \$Arg6.NbVal)\\
\end{tabarg}

\medskip
%%%%%SORTIES%%%%
\begin{tabarg}
\multicolumn{3}{l}{\titretab {Sorties}}\\
\hline
$X_{out}$ & Vecteur, \texttt{double} & $X_{in}$ moyenn� et �chantillonn� sur
la base de $time_{out}$ (m�me unit� que $X_{in})$\\
\end{tabarg}


%%%%%FORMULE%%%%
\itemstyle{Formule}

%% Formule de l'operateur : MoyEch_tps_impos_dbl
%% Dans ce fichier, les �ventuels accents sont cod�s
%% en ISO-8859-15 (aka latin9)
%% La ligne suivante donne trois exemples 
%% e accent aigu : �, e accent grave : �, e-dans-l'o : �
\[
X_{out}[i] = \frac{
\displaystyle{\sum_{t[i] - ecart\_avant}^{t[i] + ecart\_apres} X_{in}[j]} }
{N} 
\qquad i=0,Nbre\_Echant(time_{out})
\]


%%%%%AUTEUR%%%%
\itemstyle{Auteur~:}
CNRM/GMEI/TRAMM

%%%%%REFERENCE%%%%
\itemstyle{R�f�rences~:}
